\chapter{Proof for the simplified Kalman gain formula}
\label{appendix:woodbury}

This appendix presents a proof for the equivalence between \cref{eq:gain-original} and \cref{eq:gain-woodbury}.

The Sherman-Morrison-Woodbury formula can be used to do away with some inverses. 
%
Given an invertible matrix $\mathbf{S}$ and compatible matrices $\mathbf{U}$ and $\mathbf{V}^\top$, the formula states that 
%
\begin{equation}
    (\mathbf{S} + \mathbf{U}\mathbf{V}^\top)^{-1}
    =
    \mathbf{S}^{-1} -
    \mathbf{S}^{-1}\mathbf{U} 
    (\mathbf{I} + \mathbf{V}^\top\mathbf{S}^{-1}\mathbf{U})^{-1}
    \mathbf{V}^\top\mathbf{S}^{-1}.
\label{eq:woodbury}
\end{equation}

Substituting $\mathbf{V}=\mathbf{H}$, $\mathbf{U}=\mathbf{H}$ and $\mathbf{S} = \mathbf{C}^1_e$, 
\cref{eq:woodbury} becomes
%
\begin{align}
    (\mathbf{C}^1_e)^{-1} -
    (\mathbf{C}^1_e)^{-1}\mathbf{H} 
    (\mathbf{I} + \mathbf{H}^\top(\mathbf{C}^1_e)^{-1}\mathbf{H})^{-1}
    \mathbf{H}^\top(\mathbf{C}^1_e)^{-1}
    =
    ((\mathbf{C}^1_e) + \mathbf{H}\mathbf{H}^\top)^{-1}
\end{align}